\section{\sc ARL Organizational Impact} 


Dr.~Tschopp has significantly influenced the Army Research Laboratory through leadership at both the regional and branch levels. His efforts have shaped strategy, structure, and external engagement across ARL.










As ARL Central’s inaugural \textbf{Regional Lead}, Dr.~Tschopp defined and operationalized the Army Research Laboratory’s presence in the Midwest, transforming a nascent concept into a thriving regional engine for innovation, partnerships, and research impact. He led ARL’s strategic adaptation during the COVID-19 pandemic—transitioning its hub from the University of Chicago to downtown Chicago’s Discovery Partners Institute—to better align with emerging DoD priorities and help launch the Midwest Defense Innovation Hub. Under his leadership, ARL Central expanded to include researchers and major initiatives in propulsion and power, advanced materials and manufacturing, energetic materials and impact physics, the Internet of Battlefield Things (IoBT), cyber-enabled electromagnetic spectrum, and artificial intelligence/machine learning. Dr.~Tschopp built coalitions with DIU, NSIN, the 75th Innovation Command, and a broad network of academic and industry collaborators across the Midwest. Over five years, he developed and executed ARL Central’s strategic vision, guided organizational change, mentored junior scientists, launched pilot programs, and drove culture change as the inaugural Chair of ARL’s 30-member Workforce Enrichment, Recruitment, and Retention Council. He also continued to lead technically as an ARL Fellow and nationally recognized scientist, strengthening internal and external partnerships to accelerate ARL’s mission. The results are measurable: up to \$40M/year in new research investments, expanded congressional engagement, the formal establishment of the ARL Center for UAS Propulsion, the first collaborative research alliance led from the Midwest, and a growing professional network supported by high-visibility ARL events and outreach.

As \textbf{branch chief} for the Metals Branch within the Materials \& Manufacturing Science Division, Dr.~Tschopp led a strategic restructuring effort to better align the branch’s research activities with Army priorities and customer needs. He established two complementary teams: one focused on long-term, fundamental research and another dedicated to near-term, applied research with clear relevance to Army engineering centers. This restructuring broadened the branch's material focus beyond traditional basic research topics, such as nanocrystalline materials, to include critical structural materials---specifically steels and aluminum alloys---directly tied to Army Modernization Priorities. This shift not only positioned the branch to engage more effectively with Army customers but also enabled a major collaborative initiative with external partners to develop advanced steels for combat vehicle and automotive applications.  In this leadership role, Dr.~Tschopp was responsible for financial oversight, annual planning, end-of-year performance ratings, and the supervision of more than two dozen employees. His exceptional performance as branch chief was recognized with a laboratory-level award.


